% TOP_PROBLEM: {"id": 4800019, "title": "Multiple ergodic averages — Problem 19", "category_id": 11, "category": {"id": 11, "name": "dynamical_systems", "display_name": "Dynamical Systems", "description": "Problems about long-term behavior of deterministic systems, Hamiltonian dynamics, and periodic orbits.", "slug": "dynamical-systems", "order_index": 11, "created_at": "2026-07-31T15:26:25.671Z"}, "status": "partially_solved", "problem_number": "AMR-047-0019", "set_id": 13}
% FIRST_POSTED: 2026-10-01
% ENTRY_KIND: statement_only
% UnsolvedMath ID 4800019; source catalogue code AMR-047-0019
% !TeX program = lualatex
% Definition and sources only; this file is not an LLM solution attempt.
\documentclass[11pt,a4paper]{article}
\usepackage[margin=25mm]{geometry}
\usepackage{fontspec}
\setmainfont{lmroman10-regular.otf}[BoldFont=lmroman10-bold.otf,
 ItalicFont=lmroman10-italic.otf,BoldItalicFont=lmroman10-bolditalic.otf]
\usepackage{amsmath,amssymb,mathtools}
\usepackage{unicode-math}
\setmathfont{latinmodern-math.otf}
\AtBeginDocument{\let\setminus\smallsetminus}
\usepackage{xurl,hyperref}
\hypersetup{colorlinks=true,urlcolor=blue}
\setcounter{secnumdepth}{0}
\setlength{\parindent}{0pt}
\setlength{\parskip}{5pt}
\setlength{\emergencystretch}{3em}
\begin{document}
\section*{Multiple ergodic averages — Problem 19}\label{4800019}
{\small UnsolvedMath ID 4800019 · AMR-047-0019 · Dynamical Systems · AMR Open Problem Lists}
\subsection{Definitions and mathematical statement}
Let $(X,\mathcal B,\mu)$ be a probability space, and let
$T,S:X\to X$ be invertible measure-preserving transformations
which commute. Equalities between measurable transformations
may be understood modulo null sets. For bounded measurable
complex functions $f,g$, define
\[
                 A_N(x)=\frac1N\sum_{n=1}^N
                              f(T^n x)g(S^n x).
\]
Prove that $A_N(x)$ converges to a finite complex number
for $\mu$-almost every $x$, for every such system and
every $f,g\in L^\infty(\mu)$.

The exceptional set may depend on $T,S,f,g$; there
is no requirement of one common exceptional set for
all bounded measurable functions. No ergodicity or
mixing assumption is made on the individual maps or
on their joint action. In particular, $S$ is not
required to be a power of $T$.

The target concerns almost-everywhere convergence along
all averaging lengths $N$, without passage to a
subsequence. It is stronger than norm convergence
in $L^2(\mu)$ and does not follow from such norm
convergence alone. No formula for the limit is
prescribed, and the product of the separate ergodic
means is not assumed to be that limit. The
commutation hypothesis is a substantive part of the
question: the average combines the two orbits through
the same starting point and the same time index,
rather than averaging independently over two times.
\subsection{Short English statement}
Prove almost-everywhere convergence of the double average along two arbitrary commuting measure-preserving transformations.
\subsection{Sources}
\begin{itemize}
\item[S1] ULAM.AI and the UnsolvedMath Contributors, supplied problems.json, record 4800019 (AMR-047-0019), AMR Open Problem Lists. Catalogue identity and source transcription; CC BY 4.0.
\url{https://huggingface.co/datasets/ulamai/UnsolvedMath}.
\item[S2] Frantzikinakis - Some open problems on multiple ergodic averages (2016), Problem 19. Original source indexed by AMR-047-0019.
\url{https://arxiv.org/abs/1103.3808}.
\end{itemize}
\end{document}
